\documentclass[10pt,a4paper]{amsart}
\usepackage[margin=19mm]{geometry}
\usepackage{amsmath,amssymb,mathtools}
\usepackage[scr=rsfs,cal=euler]{mathalfa}
\usepackage{xfrac}
\usepackage[unicode,hidelinks,bookmarks=false]{hyperref}
\newcommand{\ie}{i.e.\ }
\newcommand{\cf}{cf.\ }
\newcommand{\resp}{respectively\ }
\newcommand{\Subsection}{\subsection}
\providecommand{\nobreakdash}{\nobreak\hbox{-}}
\setlength{\emergencystretch}{2em}
\multlinegap0pt
\usepackage{amssymb}
\newtheorem{theorem}{Theorem}
\title{Common Fixed Points of Co-Cyclic and Sub-Cyclic Mappings in Complete Metric Spaces}
\author{Dev Raj Joshi, Sanjay Mishra, Piyush Kumar Tripathi}
\date{Source: arXiv:2609.12492v1 (CC BY 4.0)}
\begin{document}
\maketitle

\noindent\emph{Source and licence.} Common Fixed Points of Co-Cyclic and Sub-Cyclic Mappings in Complete Metric Spaces. Version arXiv:2609.12492v1 (CC BY 4.0), distributed by arXiv under the Creative Commons Attribution 4.0 International licence, http://creativecommons.org/licenses/by/4.0/, as recorded in the arXiv licence metadata for this exact version and shown on the abstract page https://arxiv.org/abs/2609.12492v1. Full licence notice: \texttt{licenses/arxiv-2609.12492-CC-BY-4.0.txt}.\par\smallskip

\noindent\textit{Editorial scope.} Counted unit: the article's theorem thm:sub-cyclic-common-fixed-point (source0.tex lines 502--520). The article's complete original proof of it is delivered with it (source0.tex lines 522--663, one \texttt{proof} environment of 142 lines). No mathematics is written, repaired or re-derived: the statement and the proof are the article's own lines, in the article's own notation, and no retained line is substituted or re-flowed.  No further source-local context is needed: every cross-reference of every retained line resolves inside the two delivered spans.  Nothing was omitted from any retained span, so the package carries no omission record.  No retained line contains a citation, so the wrapper carries no bibliography and no bibliography entry.  No retained line contains a figure environment, an image-inclusion command or drawing code, so no image is delivered and the package declares no asset dependency.  Editorial formatting only, in addition: an AMS \texttt{amsart} preamble that restates the article's own declarations and macros used by the retained lines, namely the environments \texttt{theorem} on separate counters, together with the standard packages those lines need.  The retained environments print in this document as Theorem 1 in order of appearance, whereas the article prints the same items with the numbering of its own full document.  The complete native source of the article as distributed in the arXiv e-print for this exact version is bundled unchanged as \texttt{sources/210119/source0.tex}.
\begin{theorem}\label{thm:sub-cyclic-common-fixed-point}
Let $V$ and $W$ be two non-empty closed subsets of a complete metric
space $(M,\mu)$. Let $f,g:V\cup W\to V\cup W$ be two continuous self-mappings such that:

\begin{enumerate}
    \item[(i)] $f$ is sub-cyclic in $g$, that is, $ f(V)\subseteq g(W)\subseteq W$ and $ f(W)\subseteq g(V)\subseteq V$;
   

    \item[(ii)] there exists $q\in(0,1)$ such that
    \[
    \mu(f(x),f(y))
    \leq
    q\,\mu(g(x),g(y)),
    \qquad \forall x\in V,\ \forall y\in W.
    \]
\end{enumerate}
If $f$ and $g$ are compatible, then they have a unique common fixed
point in $V\cap W$.
\end{theorem}

\begin{proof}
Choose an arbitrary point $x_0\in V$. Since $f(V)\subseteq g(W)$, 
there exists $x_1\in W$ such that $f(x_0)=g(x_1)$. Again, since $f(W)\subseteq g(V)$, 
there exists $x_2\in V$ such that $f(x_1)=g(x_2)$. 

Continuing inductively, we obtain a sequence $\{x_n\}$ in $V\cup W$
such that
\[
x_{2n}\in V,\qquad x_{2n+1}\in W,
\]
and
\[
f(x_n)=g(x_{n+1}),
\qquad n\geq0.
\]

Since $x_n$ and $x_{n+1}$ belong alternately to $V$ and $W$, the
contractive condition gives
\[
\begin{aligned}
\mu(g(x_{n+1}),g(x_{n+2}))
&=
\mu(f(x_n),f(x_{n+1}))\\
&\leq
q\,\mu(g(x_n),g(x_{n+1})).
\end{aligned}
\]
Therefore, by induction,
\[
\mu(g(x_n),g(x_{n+1}))
\leq
q^n\mu(g(x_0),g(x_1)),
\qquad n\geq0.
\]

We now show that $\{g(x_n)\}$ is a Cauchy sequence. Let $m>n$. By the
triangle inequality,
\[
\begin{aligned}
\mu(g(x_n),g(x_m))
&\leq
\sum_{k=n}^{m-1}
\mu(g(x_k),g(x_{k+1}))\\
&\leq
\sum_{k=n}^{m-1}
q^k\mu(g(x_0),g(x_1))\\
&\leq
\frac{q^n}{1-q}\mu(g(x_0),g(x_1)).
\end{aligned}
\]
Since $q\in(0,1)$, the right-hand side tends to zero as $n\to\infty$.
Thus $\{g(x_n)\}$ is a Cauchy sequence.

Since $(M,\mu)$ is complete, there exists $z\in M$ such that $g(x_n)\longrightarrow z$. Moreover,
\[
f(x_n)=g(x_{n+1})\longrightarrow z.
\]
Hence
\[
f(x_n)\longrightarrow z
\quad\text{and}\quad
g(x_n)\longrightarrow z.
\]

Now observe that
\[
g(x_{2n})=f(x_{2n-1})\in f(W)\subseteq V
\]
and
\[
g(x_{2n+1})=f(x_{2n})\in f(V)\subseteq W.
\]
Thus the even subsequence $\{g(x_{2n})\}$ lies in $V$, while the odd
subsequence $\{g(x_{2n+1})\}$ lies in $W$. Since both subsequences
converge to $z$ and $V$ and $W$ are closed, we obtain $z\in V\cap W$. 

Since $f$ and $g$ are compatible and
\[
f(x_n)\longrightarrow z
\quad\text{and}\quad
g(x_n)\longrightarrow z,
\]
we have
\[
\mu(fg(x_n),gf(x_n))\longrightarrow0.
\]
By the continuity of $f$ and $g$,
\[
fg(x_n)=f(g(x_n))\longrightarrow f(z)
\]
and
\[
gf(x_n)=g(f(x_n))\longrightarrow g(z).
\]
Therefore, $\mu(f(z),g(z))=0$, 
and hence $f(z)=g(z)$. 

We now show that $z$ is a common fixed point. Since $z\in V\cap W$, we
may apply the contractive condition to $z$ and $x_n$ for each $n$ by
choosing the appropriate order of the two points. Thus,
\[
\mu(f(z),f(x_n))
\leq
q\,\mu(g(z),g(x_n)).
\]
Taking the limit as $n\to\infty$, we obtain
\[
\mu(f(z),z)
\leq
q\,\mu(g(z),z).
\]
Since $f(z)=g(z)$, it follows that
\[
\mu(g(z),z)
\leq
q\,\mu(g(z),z).
\]
Therefore, $(1-q)\mu(g(z),z)\leq0$. 
Since $q\in(0,1)$, we obtain $g(z)=z$. 
Consequently, $f(z)=g(z)=z$. Hence $z$ is a common fixed point of $f$ and $g$.

Finally, suppose that $u,v\in V\cap W$ are two common fixed points of
$f$ and $g$. Then
\[
f(u)=g(u)=u
\quad\text{and}\quad
f(v)=g(v)=v.
\]
Applying the contractive condition, we obtain
\[
\mu(u,v)
=
\mu(f(u),f(v))
\leq
q\,\mu(g(u),g(v))
=
q\,\mu(u,v).
\]
Thus, $(1-q)\mu(u,v)\leq0$. 
Since $1-q>0$, we obtain $\mu(u,v)=0$, 
and hence $u=v$. Therefore, the common fixed point is unique.
\end{proof}

\end{document}
